\documentclass[10pt,a4paper]{article}
\usepackage[margin=0.5in]{geometry}
\usepackage{amsmath}
\usepackage{graphicx}
\usepackage{booktabs}
\usepackage{setspace}
\begin{document}

\title{\vspace{-10mm}\textbf{Quarterly Technical Report: System Scalability}}
\author{Engineering Architecture Team}
\date{October 2026}
\maketitle

\section{Executive Summary}
During the third quarter, the platform experienced a 300\% increase in traffic. This technical report details the infrastructure changes implemented to ensure zero downtime and maintain sub-100ms P99 latency. The legacy monolith was completely decoupled. All core services now operate independently, communicating via asynchronous message queues to prevent cascading failures.

\section{Architecture Diagram}
\begin{figure}[h!]
\begin{center}
\includegraphics[width=0.4\linewidth]{figure_001.png}
\end{center}
\vspace{-4mm}
\caption{System diagram showing the unidirectional flow.}
\end{figure}

\section{Deployment Strategy}
The rollout was executed in three phases to mitigate risk. Phase one involved deploying the new services to a staging environment and mirroring live traffic. This allowed us to identify bottlenecks in the message queuing system before they impacted actual end users. 

Phase two was a canary release, routing 5\% of user traffic to the new architecture. After monitoring error rates and latency for 48 hours, we proceeded to phase three, which consisted of a full cutover. The entire process was automated using our internal CI/CD pipelines, allowing for instantaneous rollback if any critical alerts were triggered.

\section{Cost Analysis}
By moving to a containerized microservice approach, we were able to significantly optimize our resource utilization. Previously, the monolithic application required over-provisioning of large compute instances to handle peak loads. Now, horizontal pod autoscaling allows us to dynamically allocate resources based on real-time demand. This has resulted in a 35\% reduction in our monthly cloud infrastructure bill, despite the substantial increase in overall traffic.

\section{Conclusion}
The migration was completed successfully within the allocated budget. Future work will explore multi-region failover.
\end{document}
